% Chapter 2.
\chapter{Ageing in the field}
\label{ch:data}

\section{Eleven sites, three years}
\label{sec:sites}

The data come from eleven storage sites commissioned between 2020 and
2022 in the Netherlands and Belgium (\cref{tab:sites}). All provide
frequency containment reserve, and four also trade on the intraday
market in the evenings. Seven use lithium iron phosphate cells and four
use a nickel manganese cobalt chemistry; capacities range from 6 to
48\,MWh. Each site logs pack voltage, current and the temperature of
every module at one-second resolution, and each underwent between four
and seven reference performance tests during the observation period.

\begin{table}[tb]
  \caption[The eleven storage sites]{The eleven storage sites. Capacity
    is the nameplate value at commissioning; fade is the loss in reference
    capacity per year of operation.}
  \label{tab:sites}
  \centering
  \begin{tabular}{@{}l l r r r r@{}}
    \toprule
    Site & Chemistry & Capacity & Mean SoC & Mean temp. & Fade \\
         &           & (MWh)    &          & (°C)       & (\%/yr) \\
    \midrule
    NL-01 & LFP &  12.0 & 0.51 & 24.8 & 1.62 \\
    NL-02 & LFP &  20.0 & 0.48 & 23.1 & 1.41 \\
    NL-03 & NMC &  48.0 & 0.55 & 26.4 & 2.87 \\
    NL-04 & LFP &   6.0 & 0.62 & 25.9 & 2.15 \\
    NL-05 & NMC &  30.0 & 0.50 & 27.2 & 2.64 \\
    BE-01 & LFP &  16.0 & 0.47 & 22.6 & 1.33 \\
    BE-02 & LFP &  18.0 & 0.58 & 24.0 & 1.88 \\
    BE-03 & NMC &  24.0 & 0.66 & 28.1 & 3.52 \\
    BE-04 & LFP &  10.0 & 0.50 & 23.7 & 1.49 \\
    BE-05 & NMC &  18.0 & 0.52 & 25.5 & 2.41 \\
    BE-06 & LFP &  12.0 & 0.45 & 21.9 & 1.21 \\
    \midrule
    Total &     & 214.0 &      &      &      \\
    \bottomrule
  \end{tabular}
\end{table}

Reference capacity tests are the ground truth of this thesis. They were
performed at a constant current of C/5 between the manufacturer's voltage
limits, after a rest of at least four hours at 25\,°C ambient
(\cref{app:protocol}). Because the operators scheduled them around
market commitments, the tests are unevenly spaced, and two sites missed
one test each because of inverter faults.

\section{What drives the fade}
\label{sec:drivers}

\Cref{fig:fade} shows the reference capacity of three representative
sites. Two features stand out. First, fade is far from linear in
cycles: site BE-03, which cycles the least of the three, ages the
fastest. Second, the curves bend downward in their first year and then
straighten, the signature of a process that slows as the square root of
time.

\begin{figure}[tb]
  \marginfigure{%
  \begin{tikzpicture}
    \begin{axis}[
        width=0.98\linewidth, height=0.62\linewidth,
        xmin=0, xmax=36, ymin=0.885, ymax=1.005,
        xtick={0,6,...,36}, ytick={0.90,0.925,0.95,0.975,1.0},
        yticklabel style={/pgf/number format/.cd, fixed, precision=3},
        xlabel={Months in service}, ylabel={State of health},
        legend pos=south west,
      ]
      \addplot[accent, line width=1pt, domain=0:36, samples=60]
        {1 - 0.0118*sqrt(x) - 0.00055*x};
      \addlegendentry{BE-06 (LFP, low SoC)}
      \addplot[accent!55, line width=1pt, domain=0:36, samples=60]
        {1 - 0.0150*sqrt(x) - 0.00095*x};
      \addlegendentry{NL-04 (LFP, high SoC)}
      \addplot[black!70, line width=1pt, dashed, domain=0:36, samples=60]
        {1 - 0.0205*sqrt(x) - 0.0016*x};
      \addlegendentry{BE-03 (NMC, high SoC)}
      \addplot[only marks, mark=*, mark size=1.4pt, accent]
        coordinates {(0,1) (7,0.9653) (13,0.9506) (20,0.9356) (27,0.9241) (35,0.9108)};
      \addplot[only marks, mark=*, mark size=1.4pt, accent!55]
        coordinates {(0,1) (6,0.9578) (12,0.9367) (19,0.9177) (25,0.9017) (33,0.8826)};
      \addplot[only marks, mark=o, mark size=1.6pt, black!70]
        coordinates {(0,1) (6,0.9395) (14,0.9040) (21,0.8728)};
    \end{axis}
  \end{tikzpicture}}{%
  \caption[Reference capacity of three sites]{Reference capacity of
    three sites. Markers are reference tests; lines are the fitted ageing
    law of \cref{eq:ageing}.}
  \label{fig:fade}}
\end{figure}

To separate the drivers, the capacity loss of each site is modelled as
the sum of a calendar term, which depends on time, temperature and state
of charge, and a cycling term, which depends on charge throughput:
\begin{equation}
  1 - \mathrm{SoH}(t) =
    \underbrace{k_\mathrm{cal}\,
      e^{-E_a / R T}\, g(\sigma)\, \sqrt{t}}_{\text{calendar}}
    + \underbrace{k_\mathrm{cyc}\, \Phi(t)}_{\text{cycling}},
  \label{eq:ageing}
\end{equation}
where \(T\) is the mean cell temperature, \(\sigma\) the mean state of
charge, \(g\) an increasing function of it and \(\Phi\) the cumulative
charge throughput in equivalent full cycles. The activation energy
\(E_a\) is fixed at its laboratory value; the two rate constants are fitted
per chemistry across all sites.

\begin{keyresult}{Calendar ageing dominates}{calendar}
  Across all eleven sites, the calendar term of \cref{eq:ageing}
  accounts for 78\,\% of the capacity lost after three years (range
  64--91\,\%). A site's mean state of charge predicts its fade rate better
  than its cycle count, with a rank correlation of 0.84 against 0.21.
\end{keyresult}

The practical consequence is that the equivalent full cycle, the unit in
which warranties are written, measures the wrong thing for this kind of
duty. A site that idles at 65\,\% charge in a warm container will exhaust
its warranty in calendar years long before it reaches its cycle limit.

\subsection{Temperature inside the rack}
\label{sec:temperature}

Mean temperature hides a gradient. Within a rack, modules near the
cooling inlet run up to six degrees cooler than those at the top, and the
reference test measures only the pack, whose capacity is limited by its
weakest module. Using the hottest module temperature in place of the
mean improves the fit of \cref{eq:ageing} on every site, reducing the
residual standard deviation by 18\,\% on average. The remaining residual
is structured: it correlates with firmware changes and with seasons of
heavy intraday trading, which motivates the data-driven correction of the
next chapter.
